\documentclass[12pt,twoside]{report}

\newcommand{\reporttitle}{Type-safe Webservices Generation: Interim Report}
\newcommand{\reportauthor}{Anson Miu}
\newcommand{\reporttype}{MEng Individual Project}
\newcommand{\supervisor}{Prof. Nobuko Yoshida}
\newcommand{\sndmarker}{TBC}
\newcommand{\cid}{your college-id number}

% include files that load packages and define macros
%%%%%%%%%%%%%%%%%%%%%%%%%%%%%%%%%%%%%%%%%
% University Assignment Title Page 
% LaTeX Template
% Version 1.0 (27/12/12)
%
% This template has been downloaded from:
% http://www.LaTeXTemplates.com
%
% Original author:
% WikiBooks (http://en.wikibooks.org/wiki/LaTeX/Title_Creation)
%
% License:
% CC BY-NC-SA 3.0 (http://creativecommons.org/licenses/by-nc-sa/3.0/)
% 
% Instructions for using this template:
% This title page is capable of being compiled as is. This is not useful for 
% including it in another document. To do this, you have two options: 
%
% 1) Copy/paste everything between \begin{document} and \end{document} 
% starting at \begin{titlepage} and paste this into another LaTeX file where you 
% want your title page.
% OR
% 2) Remove everything outside the \begin{titlepage} and \end{titlepage} and 
% move this file to the same directory as the LaTeX file you wish to add it to. 
% Then add \input{./title_page_1.tex} to your LaTeX file where you want your
% title page.
%
%----------------------------------------------------------------------------------------
%	PACKAGES AND OTHER DOCUMENT CONFIGURATIONS
%----------------------------------------------------------------------------------------
\usepackage{ifxetex}
\usepackage{textpos}
\usepackage{natbib}
\usepackage{kpfonts}
\usepackage[a4paper,hmargin=2.8cm,vmargin=2.0cm,includeheadfoot]{geometry}
\usepackage{ifxetex}
\usepackage{stackengine}
\usepackage{tabularx,longtable,multirow,subfigure,caption}%hangcaption
\usepackage{fncylab} %formatting of labels
\usepackage{fancyhdr}
\usepackage{color}
\usepackage[tight,ugly]{units}
\usepackage{url}
\usepackage{float}
\usepackage[english]{babel}
\usepackage{amsmath}
\usepackage{graphicx}
\usepackage[colorinlistoftodos]{todonotes}
\usepackage{dsfont}
\usepackage{epstopdf} % automatically replace .eps with .pdf in graphics
\usepackage{natbib}
\usepackage{backref}
\usepackage{array}
\usepackage{latexsym}
\usepackage{etoolbox}

\usepackage{enumerate} % for numbering with [a)] format 

\usepackage{multicol}
\usepackage{setspace}



\ifxetex
\usepackage{fontspec}
\setmainfont[Scale=.8]{OpenDyslexic-Regular}
\else
\usepackage[pdftex,pagebackref,hypertexnames=false,colorlinks]{hyperref} % provide links in pdf
\hypersetup{pdftitle={},
  pdfsubject={}, 
  pdfauthor={\reportauthor},
  pdfkeywords={}, 
  pdfstartview=FitH,
  pdfpagemode={UseOutlines},% None, FullScreen, UseOutlines
  bookmarksnumbered=true, bookmarksopen=true, colorlinks,
    citecolor=black,%
    filecolor=black,%
    linkcolor=black,%
    urlcolor=black}
\usepackage[all]{hypcap}
\fi

\usepackage{tcolorbox}

% various theorems
\usepackage{ntheorem}
\theoremstyle{break}
\newtheorem{lemma}{Lemma}
\newtheorem{theorem}{Theorem}
\newtheorem{remark}{Remark}
\newtheorem{definition}{Definition}
\newtheorem{proof}{Proof}

% example-environment
\newenvironment{example}[1][]
{ 
\vspace{4mm}
\noindent\makebox[\linewidth]{\rule{\hsize}{1.5pt}}
\textbf{Example #1}\\
}
{ 
\noindent\newline\makebox[\linewidth]{\rule{\hsize}{1.0pt}}
}



%\renewcommand{\rmdefault}{pplx} % Palatino
% \renewcommand{\rmdefault}{put} % Utopia

\ifxetex
\else
\renewcommand*{\rmdefault}{bch} % Charter
\renewcommand*{\ttdefault}{cmtt} % Computer Modern Typewriter
%\renewcommand*{\rmdefault}{phv} % Helvetica
%\renewcommand*{\rmdefault}{iwona} % Avant Garde
\fi

\setlength{\parindent}{1em}  % indentation of paragraph


\setlength{\headheight}{14.5pt}
\pagestyle{fancy}
\fancyfoot[ER,OL]{\thepage}%Page no. in the left on
                                %odd pages and on right on even pages
\fancyfoot[OC,EC]{\sffamily }
\renewcommand{\headrulewidth}{0.1pt}
\renewcommand{\footrulewidth}{0.1pt}
\captionsetup{margin=10pt,font=small,labelfont=bf}


%--- chapter heading

\def\@makechapterhead#1{%
  \vspace*{10\p@}%
  {\parindent \z@ \raggedright %\sffamily
        %{\Large \MakeUppercase{\@chapapp} \space \thechapter}
        %\\
        %\hrulefill
        %\par\nobreak
        %\vskip 10\p@
    \interlinepenalty\@M
    \Huge \bfseries 
    \thechapter \space\space #1\par\nobreak
    \vskip 30\p@
  }}

%---chapter heading for \chapter*  
\def\@makeschapterhead#1{%
  \vspace*{10\p@}%
  {\parindent \z@ \raggedright
    \sffamily
    \interlinepenalty\@M
    \Huge \bfseries  
    #1\par\nobreak
    \vskip 30\p@
  }}
  



% %%%%%%%%%%%%% boxit
\def\Beginboxit
   {\par
    \vbox\bgroup
	   \hrule
	   \hbox\bgroup
		  \vrule \kern1.2pt %
		  \vbox\bgroup\kern1.2pt
   }

\def\Endboxit{%
			      \kern1.2pt
		       \egroup
		  \kern1.2pt\vrule
		\egroup
	   \hrule
	 \egroup
   }	

\newenvironment{boxit}{\Beginboxit}{\Endboxit}
\newenvironment{boxit*}{\Beginboxit\hbox to\hsize{}}{\Endboxit}



\allowdisplaybreaks

\makeatletter
\newcounter{elimination@steps}
\newcolumntype{R}[1]{>{\raggedleft\arraybackslash$}p{#1}<{$}}
\def\elimination@num@rights{}
\def\elimination@num@variables{}
\def\elimination@col@width{}
\newenvironment{elimination}[4][0]
{
    \setcounter{elimination@steps}{0}
    \def\elimination@num@rights{#1}
    \def\elimination@num@variables{#2}
    \def\elimination@col@width{#3}
    \renewcommand{\arraystretch}{#4}
    \start@align\@ne\st@rredtrue\m@ne
}
{
    \endalign
    \ignorespacesafterend
}
\newcommand{\eliminationstep}[2]
{
    \ifnum\value{elimination@steps}>0\leadsto\quad\fi
    \left[
        \ifnum\elimination@num@rights>0
            \begin{array}
            {@{}*{\elimination@num@variables}{R{\elimination@col@width}}
            |@{}*{\elimination@num@rights}{R{\elimination@col@width}}}
        \else
            \begin{array}
            {@{}*{\elimination@num@variables}{R{\elimination@col@width}}}
        \fi
            #1
        \end{array}
    \right]
    & 
    \begin{array}{l}
        #2
    \end{array}
    &%                                    moved second & here
    \addtocounter{elimination@steps}{1}
}
\makeatother

%% Fast macro for column vectors
\makeatletter  
\def\colvec#1{\expandafter\colvec@i#1,,,,,,,,,\@nil}
\def\colvec@i#1,#2,#3,#4,#5,#6,#7,#8,#9\@nil{% 
  \ifx$#2$ \begin{bmatrix}#1\end{bmatrix} \else
    \ifx$#3$ \begin{bmatrix}#1\\#2\end{bmatrix} \else
      \ifx$#4$ \begin{bmatrix}#1\\#2\\#3\end{bmatrix}\else
        \ifx$#5$ \begin{bmatrix}#1\\#2\\#3\\#4\end{bmatrix}\else
          \ifx$#6$ \begin{bmatrix}#1\\#2\\#3\\#4\\#5\end{bmatrix}\else
            \ifx$#7$ \begin{bmatrix}#1\\#2\\#3\\#4\\#5\\#6\end{bmatrix}\else
              \ifx$#8$ \begin{bmatrix}#1\\#2\\#3\\#4\\#5\\#6\\#7\end{bmatrix}\else
                 \PackageError{Column Vector}{The vector you tried to write is too big, use bmatrix instead}{Try using the bmatrix environment}
              \fi
            \fi
          \fi
        \fi
      \fi
    \fi
  \fi 
}  
\makeatother

\robustify{\colvec}

%%% Local Variables: 
%%% mode: latex
%%% TeX-master: "notes"
%%% End: 
 % various packages needed for maths etc.
% quick way of adding a figure
\newcommand{\fig}[3]{
 \begin{center}
 \scalebox{#3}{\includegraphics[#2]{#1}}
 \end{center}
}

%\newcommand*{\point}[1]{\vec{\mkern0mu#1}}
\newcommand{\ci}[0]{\perp\!\!\!\!\!\perp} % conditional independence
\newcommand{\point}[1]{{#1}} % points 
\renewcommand{\vec}[1]{{\boldsymbol{{#1}}}} % vector
\newcommand{\mat}[1]{{\boldsymbol{{#1}}}} % matrix
\newcommand{\R}[0]{\mathds{R}} % real numbers
\newcommand{\Z}[0]{\mathds{Z}} % integers
\newcommand{\N}[0]{\mathds{N}} % natural numbers
\newcommand{\nat}[0]{\mathds{N}} % natural numbers
\newcommand{\Q}[0]{\mathds{Q}} % rational numbers
\ifxetex
\newcommand{\C}[0]{\mathds{C}} % complex numbers
\else
\newcommand{\C}[0]{\mathds{C}} % complex numbers
\fi
\newcommand{\tr}[0]{\text{tr}} % trace
\renewcommand{\d}[0]{\mathrm{d}} % total derivative
\newcommand{\inv}{^{-1}} % inverse
\newcommand{\id}{\mathrm{id}} % identity mapping
\renewcommand{\dim}{\mathrm{dim}} % dimension
\newcommand{\rank}[0]{\mathrm{rk}} % rank
\newcommand{\determ}[1]{\mathrm{det}(#1)} % determinant
\newcommand{\scp}[2]{\langle #1 , #2 \rangle}
\newcommand{\kernel}[0]{\mathrm{ker}} % kernel/nullspace
\newcommand{\img}[0]{\mathrm{Im}} % image
\newcommand{\idx}[1]{{(#1)}}
\DeclareMathOperator*{\diag}{diag}
\newcommand{\E}{\mathds{E}} % expectation
\newcommand{\var}{\mathds{V}} % variance
\newcommand{\gauss}[2]{\mathcal{N}\big(#1,\,#2\big)} % gaussian distribution N(.,.)
\newcommand{\gaussx}[3]{\mathcal{N}\big(#1\,|\,#2,\,#3\big)} % gaussian distribution N(.|.,.)
\newcommand{\gaussBig}[2]{\mathcal{N}\left(#1,\,#2\right)} % see above, but with brackets that adjust to the height of the arguments
\newcommand{\gaussxBig}[3]{\mathcal{N}\left(#1\,|\,#2,\,#3\right)} % see above, but with brackets that adjust to the height of the arguments
\DeclareMathOperator{\cov}{Cov} % covariance (matrix) 
\ifxetex
\renewcommand{\T}[0]{^\top} % transpose
\else
\newcommand{\T}[0]{^\top}
\fi
% matrix determinant
\newcommand{\matdet}[1]{
\left|
\begin{matrix}
#1
\end{matrix}
\right|
}



%%% various color definitions
\definecolor{darkgreen}{rgb}{0,0.6,0}

\newcommand{\blue}[1]{{\color{blue}#1}}
\newcommand{\red}[1]{{\color{red}#1}}
\newcommand{\green}[1]{{\color{darkgreen}#1}}
\newcommand{\orange}[1]{{\color{orange}#1}}
\newcommand{\magenta}[1]{{\color{magenta}#1}}
\newcommand{\cyan}[1]{{\color{cyan}#1}}


% redefine emph
\renewcommand{\emph}[1]{\blue{\bf{#1}}}

% place a colored box around a character
\gdef\colchar#1#2{%
  \tikz[baseline]{%
  \node[anchor=base,inner sep=2pt,outer sep=0pt,fill = #2!20] {#1};
    }%
}%
 % short-hand notation and macros


\usepackage[nottoc]{tocbibind}
\setcitestyle{numbers}

% Includes

% Proof trees
\usepackage{proof}
\usepackage{bussproofs}

% Figure and table counter
\usepackage{chngcntr}
\counterwithout{figure}{chapter}
\counterwithout{table}{chapter}

% Commands
\newcommand{\term}[1]{\textit{#1}}
\newcommand{\mathref}[1]{\textsection #1}

\newcommand{\rulename}[1]{[\textsc{#1}]}

\newcommand{\piin}[2]{#1(#2)}
\newcommand{\piout}[2]{\bar{#1}\langle #2 \rangle }
\newcommand{\parti}[1]{{\color{blue}\mathbf{#1}}}

\newcommand{\send}[2]{#1!\langle #2 \rangle}
\newcommand{\recv}[2]{#1?(#2)}
\newcommand{\bra}[2]{#1 \rhd \left\{#2\right\}}
\newcommand{\sel}[3]{#1 \lhd #2 \langle #3 \rangle}

\newcommand{\tysend}[2]{#1![#2]}
\newcommand{\tyrecv}[2]{#1?[#2]}
\newcommand{\tybra}[2]{#1 \& \left\{#2\right\}}
\newcommand{\tysel}[2]{#1 \oplus \left\{#2\right\}}


% Typewriter font
% \renewcommand{\ttdefault}{pcr} % selects Courier font

% Citation style
\setcitestyle{square}


%%%%%%%%%%%%%%%%%%%%%%%%%%%%

\begin{document}
% front page
% Last modification: 2016-09-29 (Marc Deisenroth)
\begin{titlepage}

\newcommand{\HRule}{\rule{\linewidth}{0.5mm}} % Defines a new command for the horizontal lines, change thickness here


%----------------------------------------------------------------------------------------
%	LOGO SECTION
%----------------------------------------------------------------------------------------

\includegraphics[width = 4cm]{./figures/imperial}\\[0.5cm] 

\begin{center} % Center remainder of the page

%----------------------------------------------------------------------------------------
%	HEADING SECTIONS
%----------------------------------------------------------------------------------------
\textsc{\LARGE \reporttype}\\[1.5cm] 
\textsc{\Large Imperial College London}\\[0.5cm] 
\textsc{\large Department of Computing}\\[0.5cm] 
%----------------------------------------------------------------------------------------
%	TITLE SECTION
%----------------------------------------------------------------------------------------

\HRule \\[0.4cm]
{ \huge \bfseries \reporttitle}\\ % Title of your document
\HRule \\[1.5cm]
\end{center}
%----------------------------------------------------------------------------------------
%	AUTHOR SECTION
%----------------------------------------------------------------------------------------

\noindent
\begin{minipage}{.5\linewidth}
\begin{flushleft} \large
\textit{Author:}\\
\reportauthor
\end{flushleft}
\end{minipage}%
\begin{minipage}{.5\linewidth}
\begin{flushright} \large
\textit{Supervisor:}\\
\supervisor
\end{flushright}
\begin{flushright} \large
\textit{Second Marker:}\\
\sndmarker
\end{flushright}
\end{minipage}

\vspace{2cm}
\makeatletter
\centering
\@date 


\vfill % Fill the rest of the page with whitespace



\makeatother


\end{titlepage}



%%%%%%%%%%%%%%%%%%%%%%%%%%%% 

%\doublespacing
\tableofcontents
%\singlespacing

%%%%%%%%%%%%%%%%%%%%%%%%%%%% Main document

\chapter{Introduction}

% Rise of distributed programs -> main challenges

% How to ensure correctness in general: type system, data types

% Type discipline for concurrent programs - behavioural types (--> session types)

% Mainstream example of distirbuted programs: web services (microservice architecture), interactive web application; 

% Objective: 

\chapter{Background}

\section{Session Types}

%\subsection{Overview}
% Motivate different ways of modelling concurrency, narrow the focus to message passing
Web applications are one of many examples of distributed systems in practice. Distributed systems are built upon the interaction between concurrent processes, which can be implemented using the two main communication abstractions in \term{shared memory} and \term{message passing}. 

Shared memory provides processes with the impression of a logical single large monolithic memory but requires programmers to understand consistency models in order to correctly reason about the consistency of shared state.

Message passing interprets the interaction between processes as the exchange of messages, and best describes the communication transports found in web applications, ranging from the stateless request-response client-server interactions via HTTP to full-duplex communication channels via the WebSocket protocol \cite{WebSocketRFC}.

% Outline relevance of process algebra and session types as the relevant typing discipline
The process algebra $\pi$-calculus introduced by Milner in \cite{Milner1999} provides a formalism of the message passing abstraction in terms of the basic building blocks of sending and receiving processes, along with inductively defined continuation processes. The composition of these primitives allow us to describe more complex communication sessions.
Session types define the typing discipline for the $\pi$-calculus and provide reliability guarantees for communication sessions; the latter addresses a key challenge when reasoning about the correctness of distributed systems. 

% Outline existing implementations of session types in practice
Many studies are done on the practical applications of session types, from developing languages providing native session type support \cite{ATS2016} to implementing session types in existing programming languages across different paradigms.
Implementations of the latter approach differ by how they leverage the design philosophy and features provided by the programming language. For example, King et al. leveraged the expressive type system of PureScript to perform static session type checking in \cite{PureScript2019}, whilst Neykova and Yoshida introduced dynamic approaches to check the conformance of Python programs with respect to session types in \cite{Python2017}.

\subsection{Asynchronous $\pi$-calculus}\label{section_async}

% Introduce pi-calculus, monadic asynchronous pi-calculus.
The $\pi$-calculus models concurrent computation, where processes can execute in parallel and communicate via shared names.
We first consider the asynchronous $\pi$-calculus introduced by Honda and Tokoro in \cite{AsyncHonda}.
Among the many flavours of the calculus which vary depending on the application domain, we outline the variant as presented in \cite{C406Lecture}. 

Figure \ref{fig:async} defines the syntax of processes in asynchronous $\pi$-calculus; the asynchrony comes from the lack of continuation in the output process.

\begin{itemize}
\item $\mathbf{0}$ is the nil process and represents inactivity.
\item $\piout{u}{v}$ is the output process that will send value $v$ on $u$.
\item $\piin{u}{x}.P$ is the input process that, upon receiving a message on $u$, will bind the message to $x$ and carry on executing $P$ under this binding.
\item $P\mid Q$ represents the parallel composition of processes executing simultaneously.
\item $!P$ represents the parallel composition of infinite instances of $P$; more specifically, $!P \equiv P \mid {!P}$.
\item $(\nu a)~P$ represents a name restriction where any occurrence of $a$ in $P$ is local and will not interfere with other names outside the scope of $P$.
\end{itemize}


\begin{figure}[!hb]
\doublespacing
\[
\begin{array}{rlr}

P,Q ::= & & \text{Processes} \\
     & \mathbf{0} & \text{Nil Process} \\
\mid & \piout{u}{v} & \text{Output} \\
\mid & \piin{u}{x}.P & \text{Input} \\     
\mid & P \mid Q & \text{Parallel Composition} \\
\mid & !P & \text{Replication} \\
\mid & (\nu a)~P & \text{Restriction} \\

u,v ::= & & \text{Identifiers} \\
     & a, b, c,~\dots & \text{Names} \\
\mid & x, y, z,~\dots & \text{Variables} \\

\end{array}
\]
\singlespacing
\caption{Syntax of Asynchronous $\pi$-calculus}
\label{fig:async}
\end{figure}

% Motivate how this models interactions between concurrent processes through the reduction rule in operational semantics.
The operational semantics model the interaction between parallel processes. Whilst \cite{C406Lecture} presents the full operational semantics, we highlight the \rulename{Comm} reduction rule which specifically models message passing:
if the parallel composition of an input process and output process share the same name, the composition reduces to the continuation of the input process, substituting the variable $x$ with the message received $v$. We omit the definitions of substitution, free variables and free names, $\alpha$-equivalence and structural congruence; the interested reader may refer to \cite{C406Lecture}.

\[
\infer[\mbox{\rulename{Comm}}]{\piout{a}{v} \mid \piin{a}{x}.P ~\longrightarrow ~ P[v/x] }{}
\]

We additionally define a process $P$ to be \term{stuck} if $P$ is not the nil process and $P$ cannot be reduced any further. For example, the process $P = \piin{a}{x}.\mathbf{0} \mid \piout{b}{v}$ is stuck as the parallel composition of an input process and an output process that do not share the same name cannot be reduced using \rulename{Comm}. In practice, a stuck process contains communications that will never be executed.

\subsection{Binary Session Types}\label{section_bst}
% Introduce synchronous pi calculus with branching and selection
A \term{binary session} is a parallel composition of two processes, each representing a distinct participant.
In the context of web applications, a binary session would describe the interactions between client and server.
Without loss of generality, a \term{session} represents the sequence of send and receive actions of a single participant.

We introduce a \term{synchronous} session calculus inspired by \cite{MPST}. Figure \ref{fig:sync} defines the syntax formally; we briefly discuss the main components and how it differs from the variant introduced in \cite{C406Lecture}:

% polyadic synchronous
% communication must come with label
% equi-recursive recursion, assume guarded
\begin{itemize}
\item \textbf{Synchronous communication}: Processes that send a message will have a continuation which will be executed upon a successful send. 
\item \textbf{Polyadic communication}: A vector of values can be communicated at once.
\item \textbf{Branching and selection}: A branching process can offer a set of branches, each defined by its own label identifier and continuation process. A selection process can select a branch by sending the corresponding label identifier alongside the payload to the branching process.

\item \textbf{Labelled messages}: a label identifier is attached to all messages; the input process in \mathref{\ref{section_async}} is generalised as a branching process that offers one branch.
\end{itemize}

The \rulename{Comm} rule in the operational semantics for this calculus exemplifies these new additions: given a binary session between distinct participants $\parti{p}$ and $\parti{q}$ where $\parti{q}$ offers a set of labelled branches, if $\parti{p}$ selects a label offered by $\parti{q}$ and sends a vector of expressions $e_1, \dots, e_n$ that evaluate\footnote{We adopt the operational semantics for expression evaluation $e \downarrow v$ as defined in \cite{C406Lecture}.} to the corresponding vector of values $v_1, \dots, v_n$, the session reduces to a session with the continuation from the selection process composed with the continuation from the selected branch of the branching process, the latter having the variables $x_1, \dots, x_n$ substituted with the vector of values $v_1, \dots, v_n$ received.

\[
\infer[\mbox{\rulename{Comm}}]{
	\parti{p} :: \sel{\parti{q}}{l}{e_1 \dots e_n}.~P \mid 
	\parti{q} :: \bra{\parti{p}}{l_i(x_1 \dots x_n):Q_i}_{i\in I} ~\longrightarrow~ \parti{p} :: P \mid \parti{q} ::Q_j[v_k/x_k]^n_{k=1}
}{
	\begin{array}{cccccc}
	\exists	j \in I. l_j = l
	&
	e_1 \downarrow v_1 
	& \dots 
	& e_n \downarrow v_n
	&
	& \parti{p} \neq \parti{q}
	\end{array}
}
\]

\begin{figure}[!hb]
\doublespacing
\[
\begin{array}{rlr}

v ::= & \underline{n}~\mid~\texttt{true}~\mid~\texttt{false} & \text{Values} \\
e, e' ::= & & \text{Expressions} \\
	 & v & \text{Values} \\
\mid	 & x & \text{Variables} \\
\mid & e + e'~\mid~e - e' & \text{Arithmetic Operators} \\
\mid & e = e'~\mid~e < e' ~\mid~e > e' & \text{Relational Operators} \\
\mid & e \wedge e'~\mid~e \vee e' ~\mid~\neg e & \text{Logical Operators} \\
\mid & e \oplus e' & \text{Non-Determinism} \\

\parti{p} ::= & \parti{Client},~\parti{Server}& \text{Participant} \\

P,Q ::= & & \text{Processes} \\
     & \mathbf{0} & \text{Nil Process} \\
\mid & \bra{\parti{p}}{l_i(x_1 \dots x_n):P_i}_{i\in I} & \text{Branching} \\
\mid & \sel{\parti{p}}{l}{e_1 \dots e_n}.~P & \text{Selection} \\
\mid & \texttt{if}~e~\texttt{then}~P~\texttt{else}~Q & \text{Conditional} \\
\mid & \mu X.~P & \text{Recursive Process} \\
\mid & X & \text{Process Variable} \\

l, l' ::= & {\tt{``str''}} & \text{Label Identifiers} \\

\mathcal{M} ::= & \parti{p} :: P~\mid~\parti{q} :: Q & \text{Binary Session} \\
\end{array}
\]

\singlespacing
\caption{Syntax of Session Calculus with Branching, Selection and Recursion}
\label{fig:sync}
\end{figure}


Additionally, the calculus introduces:

\begin{itemize}
\item \textbf{Conditionals}: If $e \downarrow \texttt{true}$, the process $\texttt{if}~e~\texttt{then}~P~\texttt{else}~Q$ reduces to $P$; if $e \downarrow \texttt{false}$, the process $\texttt{if}~e~\texttt{then}~P~\texttt{else}~Q$ reduces to $Q$.
\item \textbf{Recursion}: Following the equirecursive approach, the occurrence of the process variable $X$ in the recursive process can be expanded into the process transparently; more specifically, $\mu X. P \equiv P[(\mu X. P)/X]$.
\end{itemize}


% Introduce binary session types
\term{Session types} represent the type theory for our session calculus. We define the syntax of session types for binary sessions in figure \ref{fig:bst}. 

\begin{figure}[!hb]
\doublespacing
\[
\begin{array}{rlr}

U ::= & \text{\tt{int}}~\mid~\text{\tt{bool}} & \text{Sorts} \\

S ::= & & \text{Session Types} \\
     & \mathbf{end} & \text{Termination} \\
\mid & \tybra{\parti{p}}{l_i(U_1 \dots U_n):S_i}_{i\in I} & \text{Branching} \\
\mid & \tysel{\parti{p}}{l_i(U_1 \dots U_n):S_i}_{i\in I} & \text{Selection} \\
\mid & \mu \mathbf{t}.~S & \text{Recursive Type} \\
\mid & \mathbf{t} & \text{Type Variable} \\
\end{array}
\]

\singlespacing
\caption{Syntax of Session Types}
\label{fig:bst}
\end{figure}

We derive the type of a process with a \term{typing judgement} of the form $\Gamma \vdash P: S$, which reads, \textit{under the typing context $\Gamma$, process $P$ has session type $S$}. 

The \term{typing context} records typing assumptions used as part of the derivation: in the case of binary session types, the context maps expressions to sorts, and process variables to session types. A typing judgement is constructed in terms of inference rules defined inductively on the structure of processes and expressions.

We present the rules for \rulename{Ty-Sel} and \rulename{Ty-Bra}, the remaining rules follow from \cite{C406Lecture} and can be trivially defined as they leverage the syntactic similarities between session types and our session calculus.

\[
\infer[\mbox{\rulename{Ty-Bra}}]{
	\Gamma \vdash \bra{\parti{p}}{l_i(x_1 \dots x_n):P_i}_{i\in I} : \tybra{\parti{p}}{l_i(U_1 \dots U_n):S_i}_{i\in I}
	}{
	\begin{array}{ccc}
	\forall i \in I.
	&
	&
    \Gamma, x_1: U_1, \dots, x_n: U_n \vdash P_i : S_i		
	\end{array}
}
\]

\[
\infer[\mbox{\rulename{Ty-Sel}}]{
	\Gamma \vdash \sel{\parti{p}}{l}{e_1 \dots e_n}.P : \tysel{\parti{p}}{l(U_1 \dots U_n): S}
	}{
	\begin{array}{cccc}
	\Gamma \vdash e_1 : U_1
	&
	\dots
	&
	\Gamma \vdash e_n : U_n
	&
    \Gamma \vdash P : S
	\end{array}
}
\]

% Guarantees for well-typed sessions expressed through duality, also notion of subtyping.
The definition of stuck processes from \mathref{\ref{section_async}} motivate the discussion of communication errors that may occur during interactions among participants. We outline two of the main classes of errors:

\begin{itemize}
\item \textbf{Deadlock}: Progress cannot be made when the two participants expect to be receiving a message from each other at the same time.
\item \textbf{Communication mismatch}: Progress cannot be made when the selection process sends a message with a label identifier not offered by the branching process; likewise, the payload sent must be compatible with the sort expected by the branching process for the selected branch.
\end{itemize}

Session types ensure that well-typed binary sessions are guaranteed to be free from these communication errors through the concept of \term{duality}. Duality defines a notion of \term{compatibility} between processes: two session types are dual with respect to each other if the communication between them (i.e. pairs of sending and receiving actions) always match (i.e. with respect to the selected label and message payload type). We define $\overline{S}$ as the dual type of $S$ in Table \ref{table:dual}.

% Row height in tables
\renewcommand{\arraystretch}{1.2}
\begin{center}
\begin{tabular}{rcl}
$\overline{\mathbf{end}}$ & = & $\mathbf{end}$ \\
$\overline{\tybra{\parti{p}}{l_i(U_1 \dots U_n):S_i}_{i\in I}}$ & = & $\tysel{\parti{q}}{l_i(U_1 \dots U_n):\overline{S_i}}_{i\in I}$ \\
$\overline{\tysel{\parti{p}}{l_i(U_1 \dots U_n):S_i}_{i\in I}}$ & = & $\tybra{\parti{q}}{l_i(U_1 \dots U_n):\overline{S_i}}_{i\in I}$ \\
$\overline{\mu \mathbf{t}. S}$ & = & $\mu t. \overline{S}$ \\
$\overline{\mathbf{t}}$ & = & $\mathbf{t}$ \\ 
\end{tabular}
\captionof{table}{Duality of Binary Session Types involving participants $\parti{p}$ and $\parti{q}$}
\label{table:dual}
\end{center}
\renewcommand{\arraystretch}{1}


Consequently, a binary session is well-typed if the participating processes are typed to be dual with respect to each other: we illustrate this in \rulename{MTy}.

\[
\infer[\mbox{\rulename{MTy}}]{
	\vdash \parti{p} :: P \mid 
	\parti{q} :: Q
}{
	\begin{array}{lr}
	\cdot \vdash P: S
	&
	\cdot \vdash Q: \overline{S}
	\end{array}
}
\]

The definition of duality alone restricts the definition of well-typed binary sessions to those where the two processes are derived to be \textit{exactly} dual types of one another. Consider the pair of session types below:

\[
\begin{array}{rl}
S_{\text{Client}} &= \tysel{\parti{Server}}{\text{Succ}(\texttt{int}): \tybra{\parti{Server}}{\text{Res}(\texttt{int}): \mathbf{end}}} \\
S_{\text{Server}} &= \tybra{\parti{Client}}{\text{Succ}(\texttt{int}): \tysel{\parti{Client}}{\text{Res}(\texttt{int}): \mathbf{end}},~\text{Quit}(): \mathbf{end}}
\end{array}
\]

Whilst $\overline{S_{\text{Client}}} \neq S_{\text{Server}}$, this pair of session types is intuitively compatible as the client is selecting a branch offered by the server, where the session types for the continuations of this branch for both participants are indeed dual.

This motivates the concept of \term{subtypes}, which allows a process to be typed by its ``supertype'' when required. $\leqslant$\footnote{The $\leqslant$ operator is also an overloaded relation on sorts to express subsorting.} defines the subtyping relation: $S \leqslant S'$ reads \textit{$S$ is a subtype of $S'$}, and is defined coinductively on the structure of session types.

We present the inference rules for \rulename{Sub-Bra} and \rulename{Sub-Sel} inspired by \cite{MPST} but adapted for polyadic communication; the intuition behind subtyping and subsorting is outlined below:

\begin{itemize}
\item \textbf{Branching}: The supertype of a branching process offers a subset of the branches and expects more specific types of payload; intuitively, if a process expects to receive an \texttt{int}, it can handle a \texttt{nat} payload.
\item \textbf{Selection}: The supertype of a selection process offers a superset of the internal choices and can send more generic types of payload; intuitively, if a process sends a \texttt{nat}, the payload is compatible with receivers expecting a more generic \texttt{int} payload.
\end{itemize}

\begin{prooftree}
\AxiomC{$\forall i \in I.$}
\AxiomC{$U'_1 \leqslant U_1 \dots U'_n \leqslant U_n$}
\AxiomC{$S_i \leqslant S'_i$}
\RightLabel{\mbox{\rulename{Sub-Bra}}}
\doubleLine
\TrinaryInfC{$
	\tybra{\parti{p}}{l_i(U_1 \dots U_n):S_i}_{i\in I \cup J}
	\leqslant
	\tybra{\parti{p}}{l_i(U'_1 \dots U'_n):S'_i}_{i\in I}
$}
\end{prooftree}

\begin{prooftree}
\AxiomC{$\forall i \in I.$}
\AxiomC{$U_1 \leqslant U'_1 \dots U_n \leqslant U'_n$}
\AxiomC{$S_i \leqslant S'_i$}
\RightLabel{\mbox{\rulename{Sub-Sel}}}
\doubleLine
\TrinaryInfC{$
	\tysel{\parti{p}}{l_i(U_1 \dots U_n):S_i}_{i\in I}
	\leqslant
	\tysel{\parti{p}}{l_i(U'_1 \dots U'_n):S'_i}_{i\in I \cup J}
$}
\end{prooftree}


We also introduce subsumption in \rulename{Ty-Sub} to incorporate the subtyping relation into the typing judgement.

\[
\infer[\mbox{\rulename{Ty-Sub}}]{
	\Gamma \vdash P : S'
}{
	\begin{array}{lr}
	\Gamma \vdash P: S
	&
	S \leqslant S'
	\end{array}
}
\]

This allows us to construct a derivation to show that the binary session 

\[
\mathcal{M} = \parti{Client} :: P_\text{Client}~\mid~\parti{Server} :: P_\text{Server}
\]

is well-typed, assuming $P_\text{Client}$ and $P_\text{Server}$ are typed $S_\text{Client}$ and $S_\text{Server}$ respectively.

\begin{prooftree}
\AxiomC{\vdots}
\UnaryInfC{$\cdot \vdash P_\text{Client}: S_\text{Client}$}
\AxiomC{\vdots}
\UnaryInfC{$\cdot \vdash P_\text{Server}: S_\text{Server}$}
\AxiomC{\vdots}
\doubleLine
\UnaryInfC{$S_\text{Server} \leqslant \overline{S_\text{Client}}$}

%\AxiomC{$\cdot \vdash P_\text{Client}: S_\text{Client}$}
%\AxiomC{$\cdot \vdash P_\text{Server}: S_\text{Server}$}
%\AxiomC{$S_\text{Server} \leqslant \overline{S_\text{Client}}$}
\RightLabel{\rulename{Ty-Sub}}
\BinaryInfC{$\cdot \vdash P_\text{Server}: \overline{S_\text{Client}}$}
\RightLabel{\rulename{MTy}}
\BinaryInfC{$
	\vdash \parti{Client} :: P_\text{Client} \mid 
	\parti{Server} :: P_\text{Server}
$}
\end{prooftree}

\subsection{Multiparty Session Types}

% Motivate the extension from BST to MPST (example of battleship game, existing distributed protocols that involve multiple participants?)
Whilst binary session types provide communication guarantees between exactly 2 participants, distributed systems generally involve more parties in practice. This is equally relevant in interactive web applications, as motivated by the \textit{Battleships} game example in \cite{PureScript2019} where the server coordinates interactions between two players. 

Whilst there is a natural syntactical extension to our session calculus for describing multiparty sessions\footnote{We also adopt the shorthand $\mathcal{M} ::= \prod^n_{i=1} \parti{p_i} :: P_i$ as used in the literature.} as

\[
\begin{array}{rl}
\mathcal{M} ::= & \parti{p_1} :: P_1~\mid~\parti{p_2} :: P_2~\mid~\dots~\mid \parti{p_n} :: P_n
\end{array}
\]

the same cannot be said for the binary session typing discipline, particularly with respect to duality. The same notion of duality does not extend to the decomposition of multiparty interactions into multiple binary sessions:  \cite{C406Lecture} and \cite{MPST} both present counterexamples of well-typed binary sessions that, when composed to represent a multiparty session, results in communication errors thus violating guarantees of well-typed sessions.

Honda et al. presents \term{multiparty session types} in \cite{MPAST} to extend the binary session typing discipline for sessions involving more than 2 participants, whilst redefining the notion of compatibility in this multiparty context. Multiparty session types are defined in terms a \term{global type}, which provides a bird's eye view of the communication protocol describing the interactions between pairs of participants. Figure \ref{fig:mpst} defines the syntax of global types inspired by \cite{MPST} and adapted to be compatible with our session calculus presented in \mathref{\ref{section_bst}}.

% Introduce global type, projection
\begin{figure}[!hb]
\doublespacing
\[
\begin{array}{rlr}

G ::= & & \text{Global Types} \\
     & \mathbf{end} & \text{Termination} \\
\mid & \parti{p} \to \parti{q}:\left\{l_i(U_1 \dots U_n).~G_i\right\}_{i \in I} & \text{Message Exchange} \\
\mid & \mu \mathbf{t}.~G & \text{Recursive Type} \\
\mid & \mathbf{t} & \text{Type Variable} \\
\end{array}
\]

\singlespacing
\caption{Syntax of Global Types}
\label{fig:mpst}
\end{figure}

To check the conformance of a participant's process against the protocol specification described by the global type, we \term{project} the global type onto the participant to obtain a session type that only preserves the interactions described by the global type that pertain to the participant. Projection is defined by the $\upharpoonright$ operator, more commonly seen in literature in its infix form as $G \upharpoonright \parti{p}$ describing the projection of global type $G$ for participant $\parti{p}$. Intuitively, the projected local type of a participant describes the protocol from the viewpoint of the participant.

More formally, projection can be interpreted as a \textit{partial function} ${\upharpoonright} :: G \times \parti{p} \rightharpoonup S$, as the projection for a participant may be undefined for an ill-formed global type; \cite{C406Lecture} presents examples of where this is the case, and \cite{MPST} presents the formal definition of projection.

% How does the notion of compatibility (which used to be duality) adapt to MPST?
The notion of compatibility in multiparty session types is still captured by \rulename{MTy}, but adapted to consider the local projections for all participants as supposed to dual types in the binary case. 

\[
\infer[\mbox{\rulename{MTy}}]{
	{\displaystyle \vdash \prod^{}_{i \in I} \parti{p_i} :: P_i ~: G} 
}{
	\begin{array}{llll}
	\forall i \in I. 
	&
	\cdot \vdash P_i: G \upharpoonright \parti{p_i}
	&
	&
	\text{pt}(G) \subseteq \left\{\parti{p_i} \mid i \in I \right\}
	\end{array}
}
\]


For a multiparty session $\mathcal{M} = \prod_{i \in I}\parti{p_i} :: P_i$ to be well-typed by a global type $G$, we require: 

\begin{enumerate}
\item All participant processes $\parti{p_i} :: P_i$ to be well-typed with respect to their corresponding well-defined projection $G \upharpoonright \parti{p_i}$, and 
\item $G$ does not describe interactions with participants not defined in $\mathcal{M}$ 
\end{enumerate}

Well-typed multiparty session enjoys the following communication guarantees as outlined in \cite{GentleMPST}:

\begin{itemize}
\item \textbf{Communication safety}: The types of sent and expected messages will always match.
\item \textbf{Protocol fidelity}: The exchange of messages between processes will abide by the global protocol.
\item \textbf{Progress}: Messages sent by a process will be eventually received, and a process waiting for a message will eventually receive one; this also means there will not be any sent but unreceived messages.
\end{itemize}

This motivates an elegant, decentralised solution for checking protocol conformance in practice: once the global type for the protocol is defined, local processes can verify their implementation against their corresponding projection in isolation, independent of each other. Figure \ref{fig:mpst_workflow}\todo{Create figure} illustrates this diagrammatically.

\begin{figure}[!h]
\caption{Type checking with Multiparty Session Types}
\label{fig:mpst_workflow}
\end{figure}
 

\section{The Scribble Protocol Language}

%\subsection{Overview}
% Introduce Scribble as an implementation of MPST
Whilst session type theory represents the type language for concurrent processes, it also forms the theoretical basis of proposals introduced to implement session types for real-world application development: the Scribble language is one such proposal.

Scribble \cite{Scribble} is a platform-independent description language for the specification of message-passing protocols. The language describes the behaviour of communicating processes at a high level of abstraction: more importantly, the description is independent from implementation details in the same way that the type signature of a function declaration is decoupled from the corresponding function definition.

% Protocol specification: show parallels between Scribble protocol and MPST global type
A Scribble protocol specification describes an agreement of how participating systems, referred to as \term{roles}, interact. The protocol stipulates the sequence of structured messages exchanged between roles; each message is labelled with a name and the type of payload carried by the message.

We present an example of a Scribble protocol in Figure \ref{fig:adder_scr} adapted from \cite{Hybrid2016}. The protocol specifies an arithmetic web service offered by a server to a client, represented by roles $S$ and $C$ respectively. The client is permitted to either:

\begin{itemize}
\item Send two \texttt{int}s attached to an \texttt{Add} message, where the server will respond with an \texttt{int} in a message labelled \texttt{Res}, and the protocol recurses; or,
\item Send a \texttt{Quit} message, where the server will respond with a \texttt{Terminate} message and the protocol ends.
\end{itemize}

The platform-independent nature of Scribble can be observed from the \texttt{type} declaration on Line 1: the developer has the freedom to specify message payload formats and data types from the target language of the implementation - in this case, aliasing the built-in Java integer as \texttt{int} throughout the protocol.

\begin{figure}[!h]
\begin{lstlisting}[language=Scribble]
type <java> "java.lang.Integer" from "rt.jar" as int;

global protocol Adder(role C, role S) {
	choice at C {
		Add(int, int)	from C to S;
		Res(int)		from S to C;
		do Adder(C, S);
	} or {
		Quit()		from C to S;
		Terminate()	from S to C;	
	}
}
\end{lstlisting}
\caption{Adder Protocol in Scribble}
\label{fig:adder_scr}
\end{figure}

The simplicity of the protocol specification language reflects the design goals for Scribble, as outlined in \cite{Scribble}, to be easy to read, write and maintain, even for developers who are not accustomed to the formalities in protocol specification. Moreover, we clearly observe the parallels between the Scribble language and multiparty session type (MPST) theory, from the homomorphic mapping between Scribble roles and MPST participants to the syntactic similarities between the specification in Figure \ref{fig:adder_scr} and the global type below written in the calculus.

\[
\begin{array}{rl}
G = \mu\mathbf{t}.\parti{C}\to\parti{S}:\{
& \text{Add}(\texttt{int}, \texttt{int}): \parti{S}\to\parti{C}: \{\text{Res}(\texttt{int}): \mathbf{t}\}, \\
& \text{Quit}(): \parti{S}\to\parti{C}: \{\text{Terminate}(): \mathbf{end}\}\\
\} &
\end{array}
\]


\subsection{Endpoint Finite State Machines}
% Show parallels between Scribble protocol-to-EFSM and MPST global-to-local projections
The protocol specification language is a component of the broader Scribble project in \cite{Scribble}; through the project, Honda et al. also aims to facilitate the development of endpoint applications that conform to user-specified protocols.

A Scribble global protocol can be projected to a role to obtain the local specification of said protocol from the role's viewpoint. This is analogous to the standard algorithmic projections in multiparty session type theory: the projected local specification, or local protocol, only preserves the interactions in which the target role is involved. Any user-defined type declarations in the global protocol will also be preserved. This allows local roles, also referred to as \term{endpoints}, to verify their implementation against their local protocol for conformance, independent of other endpoints. The communication safety guarantees from MPST theory also apply here: if the implementation for each endpoint is verified against its local protocol, the multiparty system as a whole will conform to the global protocol.

% Formalise syntax and properties of EFSM derived from well-formed protocols
To facilitate the verification process, a local protocol is converted into 
an \term{endpoint finite state machine} (EFSM). An EFSM encodes the control flow of the local protocol, where an initial state is defined and each transition from some state to a successor state corresponds to a valid IO action (i.e. sending to or receiving from another role) permitted at the endpoint at that state.

\begin{figure}[!hb]
\doublespacing
\[
\begin{array}{rlr}

\text{EFSM} ::= & \mathbb{R} \times \mathbb{L} \times \mathbb{T} \times \Sigma \times \mathbb{S} \times \delta & \text{Endpoint FSM} \\

\mathbb{R} ::= & r,~r',~\dots & \text{Role Identifiers} \\

\mathbb{L} ::= & l,~l',~\dots & \text{Message Label Identifiers} \\

\mathbb{T} ::= & \texttt{int},~\texttt{bool},~T,~T',~\dots & \text{Payload Format Types} \\

\Sigma ::= & & \text{Actions} \\
     & r!l(\tilde{T}) \quad \text{where } \tilde{T} \subseteq \mathbb{T} & \text{Output} \\
\mid & r?l(\tilde{T}) \quad \text{where } \tilde{T} \subseteq \mathbb{T} & \text{Input} \\

\mathbb{S} ::= & S,~S',\dots & \text{State Identifiers} \\

\delta ::= & \mathbb{S} \times \Sigma \rightharpoonup \mathbb{S} & \text{State Transition Function} \\

\end{array}
\]
\singlespacing
\caption{Syntax for Endpoint FSM}
\end{figure}

\begin{figure}[!hb]
\doublespacing
\[
\begin{array}{rl}

\texttt{initial}_{\mathbb{R}, \mathbb{L}, \mathbb{T}, \Sigma, \mathbb{S}, \delta}(S) \iff & \nexists S' \in \mathbb{S},\alpha \in \Sigma.~\delta(S',\alpha) = S \\
\texttt{terminal}_{\mathbb{R}, \mathbb{L}, \mathbb{T}, \Sigma, \mathbb{S}, \delta}(S) \iff & \delta(S) = \emptyset \\
\texttt{output}_{\mathbb{R}, \mathbb{L}, \mathbb{T}, \Sigma, \mathbb{S}, \delta}(S) \iff & \delta(S) = \{\,\alpha \in \Sigma \mid \exists r \in \mathbb{R}, l \in \mathbb{L}, \tilde{T} \subseteq{\mathbb{T}}.~\alpha = r!l(\tilde{T})\,\} \\
\texttt{input}_{\mathbb{R}, \mathbb{L}, \mathbb{T}, \Sigma, \mathbb{S}, \delta}(S) \iff & \delta(S) = \{\,\alpha \in \Sigma \mid \exists r \in \mathbb{R}, l \in \mathbb{L}, \tilde{T} \subseteq{\mathbb{T}}.~\alpha = r?l(\tilde{T})\,\} \\
\end{array}
\]
\singlespacing
\caption{Types of EFSM States}
\end{figure}


\section{Code Generation}

\subsection{Overview}
\begin{itemize}
\item Motivate the practical implementation of session types in programming
\item What needs to be verified - I/O behaviour (introduce the term behaviour typing) and linear channel usage
\end{itemize}

\subsection{Static Session Typing}
\begin{itemize}
\item PureScript implementation that addresses the 2 criteria statically
\item \begin{itemize}
\item Relevance for web development
\item Static behavioural typing
\item Static channel linearity by construction
\item Pros - type dependencies (expressing dependent types)
\item Cons
\end{itemize}
\item Other initiatives -- creating language with first-class channel primitives (e.g. SILL)
\end{itemize}

PureScript - \cite{PS2019}

\subsection{Runtime Monitors}

\subsection{Hybrid Session Verification}
\begin{itemize}
\item How the aforementioned 2 approaches motivated hybrid session verification
\item Workflow using Scribble toolchain 
\item Static behavioural typing, dynamic channel usage runtime checks
\item Pros - Abstract I/O interfaces
\item Pros - Input futures for `complex communication patterns' motivated by SMTP example
\item Cons - ``practical compromise'' as outlined in paper
\end{itemize}

Java - \cite{HY2016}

\subsection{Comparison}
\begin{itemize}
\item Comparison table between the 3 approaches for how they deal with behavioural typing 
\end{itemize}

\section{TypeScript}

\subsection{Overview}
\begin{itemize}
\item Background
\item Aims
\item Example code snippet
\end{itemize}

\subsection{Gradual Type System}
\begin{itemize}
\item Introduce static (Java), dynamic (Python/JS)
\item Use pros and cons of the former two systems to motivate the concept of a gradual type system
\item Shwo example of unsoundness
\end{itemize}

\subsection{Language Features}
\begin{itemize}
\item Outline features that are either (a) relevant to implementing web applications, (b) differ from the features prevalent in traditional OO (Java) or functional (PureScript) languages
\end{itemize}

\subsubsection{Promises}
\textit{Relevant for modelling `state-specific input futures' from the hybrid Java paper?}

\subsubsection{Discriminated Union}
\textit{Useful for dealing with the switch case approach for branch receives?}

\subsubsection{Index Types}
\textit{With index types, you can get the compiler to check code that uses dynamic property names.}

\subsubsection{Conditional Types}
\textit{Compare this with the type dependencies achieved in the PureScript paper?}


\chapter{Project Plan}
We aim to develop a multiparty session type-safe development workflow for building interactive full-stack TypeScript applications that conform to a communication protocol. This involves encoding session types using the TypeScript language and implementing a code generation workflow to generate the encodings.

%As we target full-stack applications, we tackle server-side and browser-side targets separately:
%
%\paragraph{Server-side targets} We will analyse the approaches presented in \cite{Hybrid2016, PureScript2019, Python2017} and assess their applicability with respect to the language features and type system offered by TypeScript and general back-end web development patterns. At the point of writing, we have prototyped several approaches for server-side session type TypeScript encodings and will compare their effectiveness through implementing example protocols, such as the \textit{Adder} example in \cite{Hybrid2016} and the \textit{Battleships} example in \cite{PureScript2019}.
%
%\paragraph{Browser-side targets} We will analyse how \cite{PureScript2019} handles API generation for front-end implementations, and in particular, study their approach to behavioural typing and guaranteeing channel linearity in an event-driven environment. \cite{MVU2019} also presents an architecture for verifying binary session type conformance of fornt-end implementations through formalisms of the \textit{Model-View-Update} (MVU) architecture and \textit{model types}. We will investigate how to extend their approach for multiparty session types using TypeScript. At the point of writing, we have explored methods to express the methodology in \cite{MVU2019} using React.js for a multiplayer game.

\section{Delivery}
At the point of writing, we have analysed the existing code generation approaches presented in \cite{Hybrid2016, Scribble, Python2017} and assessed their applicability with respect to verifying the communication aspects of interactive web applications. We have also explored approaches for encoding session types into TypeScript, motivated by existing works along with similar proposals specific to web development in \cite{PureScript2019, MVU2019}.

We plan to deliver our code generation implementation incrementally, aiming to get a working version of the end-to-end workflow first, then iteratively add support for more complex session type primitives. This minimises the risks involved in the project by ensuring we have a functional deliverable at the early stages of the project. We plan to deliver this basic end-to-end workflow by the end of February.

The subsequent months will involve adding support for other primitives (i.e. selection, choice and recursion) for multiparty sessions. We plan to measure our progress by writing example protocols for these milestones and checking that our implementation supports those protocols by the end of each month.

We will also try to connect with web developers in the community to experiment with our implementation to get feedback on ways to make it more applicable and compatible with industry practices, such that we can apply their feedback on the next iterations of our deliverables.

If time permits, we may explore possible extensions of applying supporting \textit{Explicit Connection Actions} presented in \cite{FASE2017} in our workflow to support more protocols.

\section{Timetable}
We attach a preliminary timetable (Table \ref{table:timetable}) for guidance.

\begin{figure}
\centering
\begin{tabular}{l|p{0.4\textwidth}|p{0.4\textwidth}}
Month & Milestones & Deadlines \\
\hline\hline
Jan & \begin{itemize}
\item Explore approaches for TypeScript encoding
\item Investigate methods for code generation from EFSM
\end{itemize} & \begin{itemize}
\item 24th: Interim Report
\item 25th: PLACES 2020
\end{itemize} \\
\hline
Feb & \begin{itemize}
\item Develop API generation end-to-end toolchain
\item Support binary sessions with simple send/receive
\end{itemize} & \begin{itemize}
\item 14th: Project Review
\end{itemize} \\
\hline
Mar & \begin{itemize}
\item Support multiparty sessions with simple send/receive
\end{itemize} & \begin{itemize}
\item 16th-20th: Examinations
\end{itemize}\\
\hline
Apr & \begin{itemize}
\item Support binary sessions with selection, choice and recursion
\item Example: calculator service
\end{itemize} & \begin{itemize}
\item 6th-17th: Easter break
\end{itemize}\\
\hline
May & \begin{itemize}
\item Support multiparty sessions with selection, choice and recursion
\item Example: Tic Tac Toe game
\end{itemize} & \begin{itemize}
\item 15th: Health check-up
\end{itemize}\\
\hline
Jun & \begin{itemize}
\item Complete report write-up
\end{itemize} & \begin{itemize}
\item 15th: Final Report
\item 26th: Final Archive
\end{itemize}
\end{tabular}
\captionof{table}{Project Timetable.}
\label{table:timetable}
\end{figure}


\chapter{Evaluation Plan}

\section{Correctness} \dots
% Qualitatively argue how the approach achieves behavioural typing and preserves channel linearity

% May explore formalisms if time permits

\section{Productivity} \dots
% Compare with PureScript implementation to talk about developer productivity when using PureScript APIs with ConcurUI compared to TypeScript APIs when using React

% Survey - qualitative and quantitative feedback

\section{Performance} \dots

% Lines of code?
% Overhead == size of JS file ran on browser + client
% Compare with baseline (no MPST)

\clearpage

\bibliographystyle{acm}
\bibliography{interim}

\end{document}
%%% Local Variables: 
%%% mode: latex
%%% TeX-master: t
%%% End: 
